\mode*
\section{Literature review protocol}
\label{app:litprotocol}

This appendix documents, in full, the systematic phase of the literature review
(\cref{sec:method-litreview}). The aim is reproducibility: every database, query,
screening decision, and per-citation provenance record can be re-run or audited.
The systematic phase was carried out with the \texttt{scholar} tool
\parencite{scholar}, which queries several bibliographic databases at once,
records each query verbatim, and supports large-language-model (LLM) assisted
screening against an explicit set of criteria.

\subsection{Databases}

We configured seven providers: Semantic Scholar, OpenAlex, DBLP, arXiv, Web of
Science, Scopus, and IEEE~Xplore. In this run the retained papers came from
OpenAlex, arXiv, Web of Science, Scopus, and IEEE~Xplore; Semantic Scholar
returned authorisation errors throughout and DBLP contributed nothing to the
retained set (see \cref{app:litlimitations}). Because the two
computer-science-specific indexes were the unproductive ones, the systematic
phase effectively searched general-purpose indexes only --- a fact that bears
directly on the comparison in \cref{app:litcomparison}.

\subsection{Search queries}

We combined two search strategies. First, a set of keyword queries broadening the
three strings of the manual phase (\cref{databasesandwords}) across all
constructs covered by the catalogue. Second, a research-question driven search in
which the tool used an LLM to translate one research question into
provider-specific queries (respecting each provider's syntax, e.g.\ uppercase
boolean operators for OpenAlex, Web of Science, and Scopus). All queries were
issued to all seven providers.

\begin{table}[h]
  \centering
  \small
  \begin{tabular}{p{0.92\textwidth}}
    \toprule
    Keyword queries \\
    \midrule
    programming misconceptions novice \\
    misconceptions introductory programming \\
    misconceptions CS1 programming \\
    Python programming misconceptions students \\
    novice programmers errors mistakes \\
    object-oriented programming misconceptions students \\
    variable assignment misconceptions programming \\
    loop iteration misconceptions novice programming \\
    recursion misconceptions novice programming \\
    conditional if statement misconceptions programming \\
    data type misconceptions programming students \\
    parameter passing misconceptions students programming \\
    variable evaluation novice programmers misconceptions \\
    avoiding object misconceptions teaching \\
    tracing quiz novice programming misconceptions \\
    novice programming mistakes large scale compilation errors \\
    object oriented programming comprehension novices misconceptions \\
    programming misconceptions exam assessment students \\
    students misconceptions loops conditionals variables Python \\
    \bottomrule
  \end{tabular}
  \caption{Keyword queries issued to all providers in the systematic phase.}
  \label{tab:litqueries}
\end{table}

The research-question driven search used the question \enquote{What
misconceptions do novice students have in introductory programming?}, which the
tool expanded per provider; representatively, for OpenAlex, Web of Science,
Scopus, IEEE, and arXiv:
\begin{quote}\small
  (\enquote{introductory programming} OR \enquote{computer programming}) AND
  (misconceptions OR \enquote{naive conceptions} OR difficulties) AND
  (novice students OR beginners)
\end{quote}
and for DBLP the simpler \texttt{misconceptions "introductory programming"
novice students}.

\subsection{Inclusion and exclusion criteria}

The criteria refine those of the manual phase. The most important change is to
relax the requirement that the word \emph{misconceptions} appear in the
\emph{title}: a paper qualifies if it identifies concrete misconceptions in the
title \emph{or} abstract.

\begin{description}
  \item[Include] a paper that identifies, lists, or characterises concrete
    misconceptions, errors, or conceptual difficulties that novices hold about
    introductory programming constructs (variables and assignment, data types,
    conditionals, loops and iteration, recursion, functions and parameter
    passing, objects and classes, program tracing, the notional machine), in
    Python or another imperative/object-oriented language whose constructs
    transfer to Python, at the introductory level (CS1, high school, or early
    university).
  \item[Exclude] papers that merely compare which programming language is most
    efficient to teach; papers about a tool that only \emph{measures} student
    knowledge without reporting the misconceptions themselves; papers not about
    introductory programming; papers about motivation, attitudes, or retention
    that report no specific misconceptions; and AI/LLM intervention or
    tool-performance studies that do not characterise misconceptions.
\end{description}

\subsection{Screening}

The search retrieved \num{117} unique papers after de-duplication. We set the
inclusion and exclusion criteria above as the research context and screened
every paper with LLM-assisted classification, using the already-labelled papers
as worked examples. \Cref{tab:litscreening} summarises the outcome; the most
common reasons for exclusion are reported alongside.

\begin{table}[h]
  \centering
  \begin{tabular}{lr}
    \toprule
    Stage & Papers \\
    \midrule
    Identified (unique, after de-duplication) & 117 \\
    Screened against criteria                 & 117 \\
    Excluded                                  & 93 \\
    \emph{Included (kept)}                     & \emph{24} \\
    \bottomrule
  \end{tabular}
  \quad
  \begin{tabular}{lr}
    \toprule
    Most common exclusion reason & Count \\
    \midrule
    Not focused on misconceptions & 49 \\
    Intervention study            & 44 \\
    AI-tool study                 & 28 \\
    Not introductory programming  & 12 \\
    Motivation / attitudes        & 9 \\
    \bottomrule
  \end{tabular}
  \caption{Screening flow for the systematic phase (a paper may carry more than
    one exclusion reason).}
  \label{tab:litscreening}
\end{table}

\subsection{Comparison with the manual review}
\label{app:litcomparison}

We treated the manual phase's catalogue as the baseline and asked how much of it
the systematic search recovered, and how much it added. Of the manual baseline's
thirteen DOI-identifiable misconception papers, the systematic search recovered
exactly one --- the most-cited survey \parencite{MisconceptionsSurvey2017} ---
and twenty-three of the twenty-four kept papers were \emph{new} relative to the
baseline. The two methods are therefore strongly complementary rather than
redundant.

The likeliest explanation is database coverage. The manual review, working
through Google Scholar with title-targeted queries and incidental citation
following, surfaced specific (and often older) computer-science conference
papers. Those are precisely the papers a CS-specific index would rank highly ---
but the two CS-specific indexes were unproductive in our run. The systematic
search over general-purpose indexes instead surfaced a broader and more recent
set. We read this as evidence that a tool-assisted multi-database search
\emph{extends} a manual review, and that the two should be combined ---
ideally with citation snowballing, which we deferred (\cref{app:litlimitations})
--- rather than either replacing the other.

\subsection{Verification and provenance}

Following the find--verify--record discipline, every retained reference carries a
provenance comment block immediately above its entry in \texttt{misconceptions.bib}.
Each block records the exact claim the reference backs (\textsf{CLAIM}), the
reproducible query that found it (\textsf{FOUND-VIA}), why it was chosen
(\textsf{PICKED}), a verbatim passage that entails the claim (\textsf{QUOTE} with
its location), and what was read to confirm applicability (\textsf{VERIFIED}).
The tool's currency check (\texttt{scholar verify}) was run over the session to
flag retracted, corrected, or superseded papers via Crossref. The provenance
blocks are validated mechanically, so no reference may enter the catalogue
without a recorded, verbatim justification.

\subsection{Limitations and deferred steps}
\label{app:litlimitations}

Three limitations of this run should be recorded. First, Semantic Scholar
returned authorisation errors and DBLP returned nothing to the retained set, so
the two computer-science-specific indexes did not contribute; the systematic
phase effectively searched general-purpose indexes, which depresses recall of
CS-conference papers. Second, results per provider were capped, so the search is
not exhaustive. Third, we deferred citation \emph{snowballing} (backward and
forward citation following from a seed set): this is the step most likely to
recover the construct-specific papers the manual review found, and it is the
natural next extension of this protocol. The full session --- queries, retrieved
papers, and screening decisions --- is saved and can be re-exported to reproduce
the tables above.
